\chapter{Particle--mesh interpolation, smoothing and NUFFT}\label{ch:mass-assignment}
\chapterstatus{written}{2026-11-14}

The link between particles and grids: density assignment with B-spline kernels, window deconvolution and interlacing,
interpolation of grids to particles, deposits for particles on any domain, Fourier-space filters and shifts, non-uniform FFTs,
and the sheet interpolation that uses the particles' Lagrangian lattice.

\begin{modulebox}{Interpolation, Smoothing}
\nsp{FML::INTERPOLATION}: \file{Interpolation/ParticleGridInterpolation.h} (\code{particles_to_grid},
\code{interpolate_grid_to_particle_positions}, \code{particles_to_fourier_grid_interlaced},
\code{deconvolve_window_function_fourier}), \file{Interpolation/ParticleGridInterpolationGeneric.h} (any grid with the interface),
\file{Interpolation/DistributedDeposit.h} (\code{particles_to_grid_any_domain}), \file{Interpolation/SheetInterpolation.h}.
\nsp{FML::GRID}: \file{Smoothing/SmoothingFourier.h} (\code{smoothing_filter_fourier_space}, \code{shift_grid_fourier}, \ldots).\par
Tests: \file{tests/test_interpolation.cpp}, \file{tests/test_interp_pencil.cpp}, \file{tests/test_distributed_deposit.cpp},
\file{tests/test_sheet_pencil.cpp}.
\end{modulebox}
\begin{modulebox}{NUFFT}
\nsp{FML::NUFFT}: \file{NUFFT/NUFFT.h} (\code{particles_to_fourier_grid_nufft}, \code{nufft_type1}, \code{nufft_type2},
\code{compute_force_nufft_with_kernel}, \cls{PSWFKernel}, \cls{ESKernel}), \file{NUFFT/Core.h} (the standalone core, in
development).\par
Documentation: \file{NUFFT/README.md}. Tests: \file{tests/test_nufft_domain.cpp}.
\end{modulebox}

\section{Assignment kernels}\label{sec:ma-kernels}

A particle at $\bm x_p$ is spread over the cells with the B-spline kernel of order $p$, the $p$-fold convolution of the top-hat
of one cell \cite{HockneyEastwood1988}:
\begin{equation}
1+\delta(\bm x_i)=\frac{N_{\rm mesh}^D}{N_{\rm part}}\sum_p\frac{m_p}{\bar m}\prod_{j=1}^DW_p\!\left(\frac{x_{i,j}-x_{p,j}}{H}\right),\qquad
\tilde W(\bm k)=\prod_{j=1}^D\left[\mathrm{sinc}\left(\frac{\pi k_j}{2k_{\rm Ny}}\right)\right]^p,
\end{equation}
with $H$ the cell size and $k_{\rm Ny}=\pi N_{\rm mesh}/B$; $p=1,2,3,4,5$ are NGP, CIC, TSC, PCS and PQS. Masses are used if the
particle has them (relative to the mean, so their units do not matter); with a $\delta f$ weight the particles carry only the
perturbation, $\delta=N_{\rm mesh}^D/N_{\rm part}\sum_pw_pW$, and no mean is subtracted. The result has zero mean.
\code{deconvolve_window_function_fourier} divides the Fourier modes by $\tilde W$.

The kernel of order $p$ reaches \code{get_extra_slices(method)} slices to the left and right (CIC: 0 and 1, TSC: 1 and 2, \ldots;
with \code{CELLCENTERSHIFTED} the shift of the cell centres changes which side needs one). Contributions to the extra slices are
added to their owners afterwards; a task with fewer slices than the kernel needs stops with a clear error, and particles on a slab
edge are assigned correctly for any number of tasks.

\paragraph*{Threads and determinism.} With one thread the particles are added in order. With several threads the local x-range is
split into chunks of at least $2p+2$ slices so that chunks $c$ and $c+2$ never write to the same cells: the even chunks are added in
parallel, then the odd ones, each chunk in particle order (the particle indices are grouped by a counting sort; the particles are not
moved). Every cell then receives its contributions in the same order for any number of threads above one, so results are
bit-reproducible. With too few slices per task (e.g.\ $N_{\rm mesh}=512$ on 32 tasks) the code falls back to atomic additions.

\paragraph*{Particles on another domain.} \code{particles_to_grid_any_domain} deposits particles whose decomposition differs from the
grid's, e.g.\ on a Peano--Hilbert domain, as GADGET's PM step does \cite{Springel2005}: each task deposits into a dense patch covering
the cells its particles touch (in several passes if the patch would exceed a cell budget), and the overlap with each task's part of the
grid is sent to it (all-zero blocks skipped) through the chunked exchange. It gives the same field as \code{particles_to_grid} up to the
order of the sums, deterministically for any number of threads, for slab and pencil grids.

\section{Aliasing and interlacing}\label{sec:ma-interlacing}

Sampling on a grid aliases the power at $\bm k+2k_{\rm Ny}\bm n$ onto $\bm k$. Averaging the density of copies of the particles
shifted by fractions of a cell,
\begin{equation}
\delta(\bm k)=\frac1M\sum_{j=1}^Me^{i\bm k\cdot\bm s_j}\delta_{\bm s_j}(\bm k),
\end{equation}
cancels the images whose phases add to zero \cite{Jing2005,Sefusatti2016}. \code{particles_to_fourier_grid_interlaced} supports:
$M=2$ with shifts $0$ and $(\frac12,\ldots,\frac12)H$, the usual interlacing, which cancels the images with odd $\sum_in_i$;
$M=2^{D-1}$ (4 in 3D) with the half-cell corners with an even number of halves, which cancels all images except those with all $n_i$
odd; $M=n^D$ sub-lattices; and $2^m\ge2^D$ by recursive refinement along the diagonal. The window is not deconvolved and the particle
positions are never modified. Combined with a PCS kernel, interlacing makes the measured power accurate up to close to the Nyquist
frequency (\cref{ch:power-spectra}).

\section{Interpolating grids to particles}\label{sec:ma-interpolation}

\code{interpolate_grid_to_particle_positions} evaluates a grid (or a vector of grids, e.g.\ the force components) at the particle
positions with the same kernel as the assignment. Using the same kernel for both makes the force interpolation the transpose of the
assignment, which removes the self-force and conserves momentum \cite{HockneyEastwood1988}. The forces of the PM solvers are computed
in Fourier space, $\bm F(\bm k)=-i\bm k\Phi(\bm k)$, one inverse FFT per component (\cref{ch:nbody-pm}); interpolation of a linear field
is exact for CIC and higher orders. \code{convolve_grid_with_kernel} convolves a grid with a general kernel in real space.

\section{Fourier-space filters and shifts}\label{sec:ma-smoothing}

\file{SmoothingFourier.h} works on any grid type in Fourier space:
\begin{itemize}
\item low-pass filters of scale $R$: sharp-$k$ ($kR<1$), Gaussian $e^{-(kR)^2/2}$ and top-hat (the transform of a ball in 3D,
      $3(\sin x-x\cos x)/x^3$, or of a disk in 2D, $2J_1(x)/x$), or a custom $W(k)$;
\item multiplication by a function of $k$, linear combinations of grids, whitening ($\delta_{\bm k}\to\delta_{\bm k}/|\delta_{\bm k}|$),
      convolutions of two fields (via real space, three FFTs);
\item spectral shifts $f(\bm x)\to f(\bm x\pm\bm s)$ by the phase $e^{\pm i\bm k\cdot\bm s}$, exact for band-limited fields. For even
      $N_{\rm mesh}$ the Nyquist modes are multiplied by $\cos(k_{\rm Ny}s)$ so the field stays real, so a shift followed by its inverse is
      the identity for them only when $s$ is a whole number of cells.
\end{itemize}

\section{Non-uniform FFT}\label{sec:ma-nufft}

The NUFFT computes the Fourier modes of point sources exactly to a requested tolerance $\epsilon$, without window or aliasing:
\begin{equation}
\text{type 1:}\quad F(\bm k)=\mathcal N\sum_pc_pe^{-i\bm k\cdot\bm x_p},\qquad
\text{type 2:}\quad f(\bm x_p)=\sum_{\bm k}F(\bm k)e^{i\bm k\cdot\bm x_p},
\end{equation}
for all modes with $|m_i|<N_{\rm mesh}/2$ ($\bm k=2\pi\bm m$). The algorithm follows FINUFFT \cite{Barnett2019}: spread the sources with a
compact kernel of width $w$ fine cells onto a grid $\sigma$ times finer, FFT, divide by the kernel's Fourier transform and keep the modes
of the output grid (type 2 is the reverse). For $\sigma=2$, $w\approx\log_{10}(1/\epsilon)+1$, so a particle costs $w^D$ operations (343
for $\epsilon=10^{-6}$ in 3D) and the fine grid has $\sigma^D$ times more cells than the output.

Two kernels: the default prolate spheroidal wave function of order zero (\cls{PSWFKernel}), computed natively by solving the angular
PSWF eigenproblem in an even Legendre basis \cite{Osipov2013} with the width rules of FINUFFT 2.5, and the exponential of semicircle
$e^{\beta(\sqrt{1-x^2}-1)}$ (\cls{ESKernel}). A reusable spatial ordering of the particles (32 blocks per dimension, a stable counting sort)
improves cache locality for repeated transforms while the results stay in the caller's order. \code{particles_to_fourier_grid_nufft}
gives the same normalisation as \code{particles_to_grid} followed by \code{fftw_r2c}, and \code{compute_force_nufft_with_kernel} the PM
force without window or aliasing. MPI (particles in the slab or pencil of the fine grid) and OpenMP are supported. A standalone,
general complex NUFFT core of types 1--3 independent of the particle classes is in development (\file{Core.h}).

\begin{validation}[Validation: NUFFT]
In 3D with $3\times10^5$ random particles, $64^3$ modes, $\epsilon=10^{-6}$ and $\sigma=2$, the PSWF kernel agrees with FINUFFT 2.5.1 to
about $10^{-9}$ (relative maximum error) for both types; ordered PSWF was a few per cent faster than ordered ES and 10--15\% slower than
FINUFFT (a local measurement). \file{tests/test_nufft_domain.cpp}: on pencils the type 1 modes, type 2 values (both kernels) and forces
equal those on slabs.
\end{validation}

\section{Sheet interpolation}\label{sec:ma-sheet}

Particles that started on a lattice $\bm q_i$ sample the dark-matter sheet $\bm x(\bm q)=\bm q+\bm\Psi(\bm q)$, and before shell crossing
$\bm\Psi$ is smooth, so it can be Fourier-interpolated between the particles \cite{Hahn2013,Stucker2020,List2024}. The sheet interpolation
creates tracers at
\begin{equation}
\bm x=\bm q_i+\bm s+\bm\Psi(\bm q_i+\bm s),\qquad \bm s=\frac{(j_1,\ldots,j_D)}{RN_{\rm part,1D}},\quad j_d=0,\ldots,R-1,
\end{equation}
and assigns all $R^D$ of them to the grid. This removes the shot noise and the lattice pattern of the particles (with $R$ equal to the ratio of
the force grid to the particle grid there is one tracer per cell) at the cost of $R^D$ assignments. $\bm\Psi$ comes from the particles'
Lagrangian coordinates (\code{get_q}) or is supplied (e.g.\ the 1LPT field). After shell crossing $\bm\Psi$ has kinks and the interpolation
rings, so the method is best on large scales and at early times. It is implemented for slabs and pencils.

\begin{validation}[Validation: assignment and interpolation]
\file{tests/test_interpolation.cpp}: NGP interpolation of a constant gives the constant for every particle; lattice particles on slab
edges are all assigned ($\sum\delta=0$ for NGP--PQS, also with 10 tasks); interpolation of a linear field is exact for CIC--PQS; unknown
methods and too-thin slabs stop with a clear message. \file{tests/test_interp_pencil.cpp}: interlacing of every order, offset deposits and
deconvolution give the same modes on pencils and slabs. \file{tests/test_distributed_deposit.cpp}: particles on a Peano--Hilbert domain,
scattered randomly or all on one task, NGP--PQS with masses and offsets, several passes and tiny messages, pencil grids, OpenMP
determinism, against \code{particles_to_grid}. \file{tests/test_sheet_pencil.cpp}: sheet densities on pencils equal those on slabs.
\end{validation}
